\documentclass[12pt,a4paper]{article}
% Preambule commun aux comptes rendus CR_*.tex.

\usepackage{array}
\newcolumntype{C}[1]{>{\centering\arraybackslash}m{#1}}
\usepackage{multirow}
\usepackage{geometry}
\geometry{a4paper, margin=1in}
\usepackage{graphicx}
\usepackage{float}
\usepackage{tikz}
\usetikzlibrary{arrows.meta,positioning}
\usepackage{pgfgantt}
\usepackage{pdflscape}
\usepackage{enumitem}
\usepackage{booktabs}
\usepackage{longtable}
\usepackage{ragged2e}
\usepackage{tabularx}
\usepackage{xparse}
\usepackage{ltablex}
\usepackage{xltabular}
\usepackage{subcaption}

\newcolumntype{Y}{>{\RaggedRight\arraybackslash}X}

\newcommand{\StyledTableSetup}{%
  \footnotesize
  \renewcommand{\arraystretch}{1.2}%
  \setlength{\tabcolsep}{4pt}%
}


% ============================================================
% IHEDN COVER TEMPLATE (XeLaTeX)
% ============================================================
% Compilation (from project root):
%   latexmk -pdf main.tex
%   LATEXMK_SHELL_ESCAPE=0 latexmk -pdf main.tex
%
% Optional environment controls from latexmkrc:
%   LATEXMK_ENGINE=xelatex        (default/enforced)
%   LATEXMK_SHELL_ESCAPE=1|0      (default 1, needed for SVG conversion)
%
% Required package for this template:
%   \usepackage{ihedn-cover}
%
% Note: ihedn-cover already loads needed internal packages
% (fontspec, graphicx, tikz, ragged2e, optional svg).
% ============================================================

% ----------------------------------------------------------------
% Package options (documented)
% ----------------------------------------------------------------
% Geometry scope:
%   apply-geometry            -> force A4 + zero margins while cover is built
%   no-apply-geometry         -> do not alter host geometry (default)
%
% Typesetting freeze:
%   global-typesetting        -> apply freeze globally (intrusive)
%   no-global-typesetting     -> no global freeze (default)
%
% Small-caps handling:
%   local-smallcaps-fix       -> \textsc -> \MakeUppercase only inside cover (default)
%   no-local-smallcaps-fix    -> disable local fix
%   global-smallcaps-fix      -> global override (intrusive)
%   no-global-smallcaps-fix   -> disable global fix (default)
%
% Main font handling:
%   global-mainfont           -> set Times New Roman globally (default)
%   local-mainfont            -> set Times only during cover rendering
%   no-mainfont-override      -> keep host document main font
%
% SVG support:
%   svg-support               -> enable SVG logos (default, needs shell-escape)
%   no-svg-support            -> disable SVG path handling (pdf/png only)
%
% Debug overlay:
%   debug                     -> show mm grid + bounding boxes on cover
%   no-debug                  -> hide debug overlay (default)
% ----------------------------------------------------------------
\usepackage[
  no-apply-geometry,
  no-global-typesetting,
  global-smallcaps-fix,
  local-smallcaps-fix,
  global-mainfont,
  svg-support,
  no-debug
]{ihedn-cover}

% ----------------------------------------------------------------
% Bibliography stack (BibLaTeX/Biber, style from subtree vendor/)
% ----------------------------------------------------------------
\usepackage{polyglossia}
\setdefaultlanguage{french}
\makeatletter
\g@addto@macro\captionsfrench{%
  \renewcommand{\tablename}{Tableau}%
  \renewcommand{\figurename}{Figure}%
}
\makeatother
\usepackage{csquotes}
\usepackage{xurl}
\usepackage{hyperref}
\makeatletter
\let\ihedn@orig@input@path\input@path
\def\input@path{{vendor/biblatex-sciences-po-fr/style/}}
\makeatother
\usepackage[
  backend=biber,
  style=sciencespo,
  hyperref=true,
  autolang=other
]{biblatex}
\AtBeginBibliography{\sloppy\setlength{\emergencystretch}{3em}}
\makeatletter
\let\input@path\ihedn@orig@input@path
\makeatother
\addbibresource{bibliographies/bib/rapport.bib}

% ----------------------------------------------------------------
% tcolorbox
% ----------------------------------------------------------------
\usepackage[most]{tcolorbox}
\tcbuselibrary{breakable}

% Example: use Arial for the full report + cover.
% Uncomment and switch package option to no-mainfont-override.
% \setmainfont{Arial}[
%   UprightFont    = *,
%   BoldFont       = * Bold,
%   ItalicFont     = * Italic,
%   BoldItalicFont = * Bold Italic
% ]

%==================================================
% Liste des propositions
%==================================================

\makeatletter

\newcommand{\listpropositionname}{Liste des propositions}

\newcommand{\listofpropositions}{
  \section*{\listpropositionname}
  \addcontentsline{toc}{section}{\listpropositionname}
  \@starttoc{lop}
}

\newcommand*\l@proposition{\@dottedtocline{1}{1.5em}{2.3em}}

\makeatother

%==================================================
% Compteur des propositions
%==================================================

\newcounter{proposition}

%==================================================
% Style graphique des propositions
%==================================================

\newtcolorbox{propositionbox}[2][]{
  lower separated=false,
  colback=white!80!gray,
  colframe=white!20!black,
  fonttitle=\bfseries,
  colbacktitle=white!30!gray,
  coltitle=black,
  enhanced,
  sharp corners,
  attach boxed title to top left={
    xshift=0.5cm,
    yshift=-2mm
  },
  title=#2,
  #1
}

%==================================================
% Environnement proposition
%==================================================

\newenvironment{proposition}[1]
{
  \refstepcounter{proposition}

  \addcontentsline{lop}{proposition}{%
    \protect\numberline{\theproposition}#1%
  }

  \begin{propositionbox}
  {Proposition~\theproposition~: #1}
}
{
  \end{propositionbox}
}

% Renommage Liste des figures en Table des figures
\addto\captionsfrench{%
  \renewcommand{\listfigurename}{Table des figures}
}

\newcommand{\heure}[2]{{\NoAutoSpacing #1:#2}}

\DeclareFieldFormat[misc]{title}{#1}

% Activer ce package avec l'option none retire la césure.
%\usepackage[none]{hyphenat}

\usepackage{adjustbox}
\usepackage{pdfpages}

\makeatletter
\newcommand{\IhednSetPdfPageCount}[2]{%
  \begingroup
  \edef\ihedn@pdfpath{#2}%
  \ifXeTeX
    \global#1=\XeTeXpdfpagecount "\ihedn@pdfpath" %
  \else
    \ifLuaTeX
      \saveimageresource{\ihedn@pdfpath}%
      \global#1=\lastsavedimageresourcepages\relax
    \else
      \pdfximage{\ihedn@pdfpath}%
      \global#1=\pdflastximagepages\relax
    \fi
  \fi
  \endgroup
}
\makeatother

\newcount\IhednIncludedPdfPageCount
\newcommand{\IhednRenderPdfPagesInBoxes}[1]{%
  \IhednSetPdfPageCount{\IhednIncludedPdfPageCount}{#1}%
  \foreach \IhednIncludedPdfPage in {1,...,\the\IhednIncludedPdfPageCount}{%
    \begin{tcolorbox}[
      colback=gray!5,
      colframe=black,
      boxrule=0.5pt,
      arc=3mm
    ]
    \centering
    \includegraphics[page=\IhednIncludedPdfPage,width=0.95\textwidth]{#1}
    \end{tcolorbox}
    \ifnum\IhednIncludedPdfPage<\IhednIncludedPdfPageCount
      \clearpage
    \fi
  }%
}

\usepackage{tablefootnote}


\begin{document}

% ----------------------------------------------------------------
% Cover keys (all available keys are shown below)
% ----------------------------------------------------------------
% Header block (top-right), logo, title, report lines, subtitle, committee.
\PageDeGardeIHEDN{
    header_1={Union-IHEDN},
    header_2={Association des auditeurs IHEDN},
    header_3={région Paris / Île de France},
    date={le 26 janvier 2026},
    logo_path={Figures/logos_IHEDN/Logo_AR_16_PARIS_IDF.jpg},
    title={Crises majeures : quel rôle pour les citoyens engagés et les associations IHEDN ?},
    title_max_lines=2,
    title_fontsize={26.000pt},
    title_baselineskip={28.667pt},
    reportline_1={Rapport du Comité d'étude de l'Association régionale des auditeurs},
    reportline_2={de l'IHEDN région Paris / Île de France (AR 16)},
    subtitle={Compte rendu de la réunion du 22 janvier 2026},
    committee={Comité d'étude, d'octobre 2025 à juin 2026, composé de : Mme Valérie Arlé-Faivre, \textit{présidente} ; M. Aurélien Duvignac-Rosa, \textit{rapporteur} ; M. Pascal Bayot Duponchel ; M. Jean-Jacques Cleynen ; Mme Valérie Cowan ; M. Alain Fontan ; M. Pierre-Marie Giard ; Mme Muriel Joyeux ; M. Bruno Leray ; M. Yves-Henri Renhas ; M. Marc Roure ; Mme Isabelle de Segonzac ; M. Gérard Turck ; Mme Marie Danielle Vasquez-Duchêne.}
}

\section{Rappel du cadre général et des travaux antérieurs}

Le comité s'inscrit dans la continuité des réflexions engagées lors des séances précédentes, ayant permis de cadrer l'approche de gestion de crises majeures dans une logique de \textbf{sécurité globale}, intégrant les dimensions naturelles, technologiques, sanitaires, informationnelles, sociétales et hybrides.
Il a également été fait mention que la désinformation constitue un \textbf{facteur transversal aggravant}, susceptible d'amplifier les effets d'une crise et d'altérer la confiance des citoyens envers les institutions.

\subsection{Rappel des quatre cercles relatifs aux notions de défense national et de sécurité nationale}

« La notion de défense nationale peut être appréhendée à travers les quatre cercles suivants qui définissent le champ d'intervention de l'IHEDN :
\begin{itemize}
\item Le premier cercle est celui de la défense militaire. Il s'agit de la défense par les armes qui se caractérisent par l'emploi ou par le perspective d'emploi de la forme armée, autrement dit de la violence physique ; ses critères sont ceux de la paix et de la guerre ;
\item Le deuxième cercle est celui de la défense nationale. Elle vise à nous protéger des menaces de tous ordres, c'est-à-dire des atteintes intentionnelles à dimension politique à notre souveraineté, nos intérêts, notre environnement ou notre cohésion nationale. Il s'agit de faire face par exemple aux menaces que constituent entre autre la cybercriminalité, le terrorisme, l'espionnage, la désinformation, le rachat d'entreprises sensibles, la territorialisation de la mer ou les menaces dans l'espace; 
\item Le troisième cercle est celui de la sécurité nationale. Il s'agit du domaine de l'anticipation, de la préparation et de la réaction aux risques. Tous les risques ne sont pas concernés, mais seulement ceux qui, par leur ampleur, concernent la communauté nationale dans son ensemble : organisation de l'État ou du système politique, infrastructures, entreprises critiques, etc. Il peut s'agir des risques naturels ou technologiques, des épidémies, de l'émergence d'une vulnérabilité ou d'une dépendance technologique ou encore de conséquences systémiques de grands trafics ;
\item Le quatrième cercle est celui de la sécurité international. La protection contre les atteintes à la défense et à la sécurité national ne saurait être garantie sans action en faveur de la sécurité internationale, qu'il s'agisse de l'action diplomatique bilatérale ou multilatérale et de la maîtrise des armements ou de la promotion des mécanismes de sécurité collective.
\end{itemize}

Pour garantir notre défense et notre sécurité nationale, c'est l'ensemble de l'État et, au-delà, toutes les forces vives de la communauté nationale, ses entreprises, des collectivités, ses forces politiques, sa jeunesse, qui doivent se mobiliser.

Il en résulte que la défense est permanente et non pas seulement une organisation du temps de guerre. Elle est globale puisqu'elle prend en compte tous les aspects militaires et non militaires de la Nation dans un cadre qui est aussi européen. Elle est enfin transverse, c'est-à-dire interministérielle et doit être pensée et mise en œuvre comme telle ».

\newpage

\subsection{Rappel des enjeux abordés relatifs à la thématique : « Crises majeures : quel rôle pour les citoyens engagés et les associations IHEDN »}

Les enjeux abordés sont les suivants :

\begin{itemize}
\item Place des réserves civiles, citoyennes et stratégiques dans la gestion des crises (pandémie, cyberattaque, catastrophe naturelle, conflit hybride, \textit{blackout} énergétique, etc.) ;
\item Mobilisation  des compétences issues des IHEDN : relais d'information stratégique, animation de la résilience locale, sensibilisation de la population , liens civilo-militaires ;
\item Rôle des associations régionales en soutien des autorités (préfectures, mairies, forces armées, sécurité civile) ;
\item Partage d'expériences locales : coordination, formation, remontée d'informations, mise en réseau  ;
\item Renforcement de la culture de sécurité globale et de l'esprit de défense élargi au niveau territorial ;
\end{itemize}

\textbf{Pertinence du thème :} Face à la montée des risques systémiques (climat, énergie, cyber, terrorisme, guerre informelle) les crises majeures sont appelées à devenir plus fréquentes. Dans ce contexte, les membres et associations IHEDN, souvent composés de cadres expérimentés issus de la société civile et des forces armées , représentent un maillon stratégique de la résilience nationale. Ce thème questionne leur capacité à agir, à fédérer et à contribuer utilement au service de l'intérêt général.

\newpage

\section{Problématique centrale du comité}

Le comité a confirmé un point fondamental : \textbf{L'IHEDN n'a pas vocation à décrire ou gérer directement une crise}, mais à \textbf{contribuer utilement à la préparation, à l'accompagnement et à l'analyse des crises}, en complément des autorités compétentes.

La réflexion porte donc sur :
\begin{itemize}
\item La \textbf{mobilisation des auditeurs IHEDN} ;
\item L'identification d'une \textbf{plus-value concrète} de la part des associations, et par extension des membres, de l'IHEDN, malgré des dispositifs de gestion de crise déjà structurés ;
\item La définition de \textbf{moments d'intervention pertinents}, notamment en amont et en aval (RETEX).

\end{itemize}

\section{Mobilisation des auditeurs et associations IHEDN}

\textbf{L'annuaire de l'IHEDN représente un vivier significatif de compétences et d'expertises, susceptible de constituer le socle d'un dispositif structuré de préparation aux crises et de retour d'expérience (RETEX).}

Sur la base de ce constat, trois enjeux ont été identifiés :
\begin{itemize}
\item Exploiter le potentiel des auditeurs IHEDN qui souhaitent s'impliquer de manière concrète ;
\item Clarifier le rôle attendu des membres dans une logique réaliste et responsable ;
\item Éviter toute confusion entre engagement citoyen, expertise et action opérationnelle.
\end{itemize}

L'ensemble de ces éléments ont permis d'aboutir à une proposition concrète de la part de notre comité : la mise en place d'un questionnaire national à destination des auditeurs IHEDN\footnote{idée émise par Monsieur Alain Fontan lors de la réunion du 11 décembre 2025, sur la base de la proposition de Monsieur Bruno Leray}. 

\section{Proposition structurante : mise en place d'un questionnaire national}

\subsection{Objectifs du questionnaire}

Le comité propose la conception d'un \textbf{questionnaire national à destination des adhérents IHEDN}, visant à :

\begin{itemize}
\item Cartographier les \textbf{domaines d'expertise} (nucléaire, fluvial, juridique, sociologique, communication, etc.) ;
\item Identifier les \textbf{ressources disponibles}, leur nature et leur ancienneté ;
\item Évaluer la \textbf{disponibilité réelle} des auditeurs ;
\item Repérer les profils volontaires susceptibles d'intégrer un dispositif d'appui ou de réserve\footnote{Une réserve IHEDN avait été constituée par le passé, sans qu'elle n'ait toutefois connu de prolongement opérationnel.}.
\end{itemize}

Les réponses permettraient de \textbf{filtrer les personnes réellement enclines à s'engager}, et de constituer une base crédible de capacités mobilisables.

\subsection{Gouvernance et calendrier}

Le calendrier de cette démarche peut se diviser de la manière suivante :
\begin{itemize}
\item Le lancement du questionnaire doit être \textbf{validé par le Bureau et le CODIR} ;
\item Une \textbf{co-construction} pourrait être envisagée avec l'IRSEM (pôle défense) ;
\item Lancement souhaité : \textbf{février}, avec clôture \textbf{fin mars} ;
\item Restitution des résultats lors de \textbf{l'Assemblée générale} qui est un moment clef pour la diffusion d'informations importantes.
\end{itemize}

\section{Structuration du questionnaire autour des trois temps de la crise}

Le comité propose d'articuler le questionnaire autour des \textbf{trois phases de la crise}.

\subsection{Avant la crise (anté-crise) : Comment préparer une crise sans générer de peur ou d'insécurité ?}

\begin{itemize}
\item Contribution à la \textbf{préparation}, à la sensibilisation et à l'animation de la résilience ;
\item Intégration éventuelle des bénévoles IHEDN dans des dispositifs existants (PCS, formations, exercices) ;
\end{itemize}

\subsection{Pendant la crise}

Comme évoqué, le rôle des associations et des membres IHEDN reste relativement modeste au sein de la crise en elle-même. \textbf{La plus-value de l'IHEDN porte notamment sur une analyse en amont et en aval de la crise.} Toutefois, deux éléments peuvent être envisagés :
\begin{itemize}
\item Rôle possible de relais d'information fiable sous une autorité préfectorale ;
\item  Clarification indispensable de la \textbf{voie hiérarchique} et du cadre d'action.
\end{itemize}

\newpage

\subsection{Après la crise : Comment analyser une crise et en tirer des enseignements utiles ?}

En aval de la crise, l'expertise issue des auditeurs IHEDN permettrait d'envisager les éléments décrits ci-après :
\begin{itemize}
\item Participation aux retours d'expérience (RETEX) ;
\item Analyse sociologique, organisationnelle et informationnelle de la crise.
\end{itemize}


\section{Apports attendus et recherche de plus-value}

\textbf{Le comité a insisté sur la nécessité d'être force de proposition, malgré des organisations de gestion de crise déjà bien établies.}

Quatre Axes ont été identifiés :

\begin{enumerate}
\item Identification des moments où l'IHEDN peut intervenir utilement ;
\item Mise en valeur des expertises rares ou transversales ;
\item Participation à des actions de formation et de préparation, notamment en matière d'évacuation ou de sensibilisation ;
\item Clarification du rôle de chacun au sein de l'IHEDN.
\end{enumerate}

Une étude du rapport de synthèse du Beauvau de la Sécurité civile peut également étoffer notre étude\footcite{MI:2025:RDS}.


\section{Relations avec les autorités et reconnaissance institutionnelle}

\subsection{Identification et légitimité}

Comme l'a souligné Madame Muriel Joyeux, \textbf{toute participation à la gestion de crise suppose une identification claire et un label reconnu, notamment par la DGSC et chaque personne engagée doit répondre à une utilité réelle et identifiée.}

\subsection{Dialogue avec les autorités}

Le comité propose \textbf{des entretiens avec des préfets, maires, responsables de la sécurité civile, pompiers}, afin de répondre à des questions simples. Notamment sur les éléments manquants aujourd'hui pour apporter des solutions pérennes. Mais également dans quelle mesure l'IHEDN et ses membres pourraient apporter des réponses concrètes. De manière plus précise encore, l'intégration des auditeurs IHEDN travaillant déjà dans les administrations concernées pourrait être envisagée.

\newpage

\section{Conclusion et axes de travail envisagés dans le cadre des prochaines réunions}

À la suite des échanges, le comité retient les actions suivantes :

\begin{itemize}
\item Identifier les associations IHEDN ayant travaillé sur des thématiques similaires et les contacter ;
\item Lancer un sondage national auprès des adhérents ;
\item Identifier, auprès des maires et des préfets, les manques et les besoins concrets ;
\item Organiser un retour structuré des résultats ;
\item Échanger avec la direction de l'IHEDN afin de clarifier la feuille de route et les marges d'action autorisées ;
\item Si le questionnaire est validé par le CODIR, utiliser cet élément pour galvaniser la mobilisation des auditeurs et renforcer la lisibilité de l'IHEDN auprès des autorités.
\end{itemize}
~\\
Le comité tient à remercier Madame Camille Guthmann, coordinatrice des référents départementaux, invitée à cette séance, pour la qualité de son intervention et le partage de son expertise, qui ont utilement nourri les réflexions du groupe.

\end{document}